\documentclass[11pt]{article}

\usepackage[margin=1in]{geometry}
\usepackage{amsmath,amssymb,amsthm}
\usepackage{mathtools}
\usepackage{enumitem}
\usepackage[hidelinks]{hyperref}

\title{Treadle Requirements for Two-Foot Shaft Selection}
\author{}
\date{\today}

\newtheorem{definition}{Definition}
\newtheorem{problem}{Problem}
\newtheorem{theorem}{Theorem}
\newtheorem{proposition}{Proposition}
\newtheorem{conjecture}{Conjecture}
\newtheorem{example}{Example}
\newtheorem{observation}{Observation}

\newcommand{\pow}{\mathcal{P}}
\newcommand{\A}{\mathcal{A}}
\newcommand{\twounions}{\langle \A\rangle_2}

\begin{document}

\maketitle

\begin{abstract}
These notes study a combinatorial design problem motivated by treadle loom
construction. If a loom has \(n\) shafts and the weaver can press at most two
treadles at once, how many treadles are required so that every nonempty
combination of shafts can be lifted? Equivalently, what is the smallest family
\(\A \subseteq \pow([n])\) such that every nonempty subset of \([n]\) is the
union of at most two members of \(\A\)?
\end{abstract}

\section{The Problem}

Let
\[
  [n] = \{1,2,\ldots,n\}
\]
denote the set of shafts. A treadle may be tied to any subset of shafts. If two
treadles are pressed simultaneously, then the lifted shafts are the union of
the two tied subsets.

\begin{definition}
  A family \(\A \subseteq \pow([n])\setminus\{\varnothing\}\) is called a
  \emph{two-foot treadle basis} for \([n]\) if every nonempty subset
  \(X \subseteq [n]\) can be written as
  \[
    X = U \cup V
  \]
  for some \(U,V \in \A\). We allow \(U=V\), which models pressing one treadle
  rather than two.
\end{definition}

\begin{definition}
  Define
  \[
    \tau_2(n)
    =
    \min\bigl\{|\A| : \A \subseteq \pow([n])\setminus\{\varnothing\}
    \text{ is a two-foot treadle basis for } [n]\bigr\}.
  \]
  Thus \(\tau_2(n)\) is the minimum number of treadles needed to make every
  nonempty shaft combination available with at most two feet.
\end{definition}

The central question is:

\begin{problem}
  Determine \(\tau_2(n)\), and describe the extremal families \(\A\) attaining
  it.
\end{problem}

\section{Small Values}

We use the shorthand \(13\) for the subset \(\{1,3\}\), \(456\) for
\(\{4,5,6\}\), and so on. The known small values are:
\[
\begin{array}{c|c|l}
  n & \tau_2(n) & \text{example of an optimal family } \A\\
  \hline
  1 & 1  & \{1\}\\
  2 & 2  & \{1,2\}\\
  3 & 4  & \{1,2,3,12\}\\
  4 & 6  & \{1,2,3,4,13,24\}\\
  5 & 10 & \{1,2,3,4,5,12,13,23,45,123\}\\
  6 & 14 & \{1,2,3,4,5,6,12,23,13,45,56,46,123,456\}
\end{array}
\]

These examples all come from the same construction: split the shafts into two
blocks and include every nonempty subset inside each block.

\section{The Split Construction}

\begin{proposition}[Split construction]
  Let \([n]=X\sqcup Y\) be a partition into two disjoint blocks. Define
  \[
    \A_{X,Y}
    =
    \bigl(\pow(X)\setminus\{\varnothing\}\bigr)
    \cup
    \bigl(\pow(Y)\setminus\{\varnothing\}\bigr).
  \]
  Then \(\A_{X,Y}\) is a two-foot treadle basis for \([n]\).
\end{proposition}

\begin{proof}
  Let \(S\subseteq[n]\) be nonempty. Then
  \[
    S = (S\cap X)\cup(S\cap Y).
  \]
  If both intersections are nonempty, each is a member of \(\A_{X,Y}\). If one
  intersection is empty, then \(S\) itself is a nonempty subset of the other
  block, hence \(S\in\A_{X,Y}\). Therefore \(S\) is the union of at most two
  members of \(\A_{X,Y}\).
\end{proof}

Choosing \(X\) and \(Y\) as evenly as possible gives
\[
  \tau_2(n)
  \le
  \left(2^{\lfloor n/2\rfloor}-1\right)
  +
  \left(2^{\lceil n/2\rceil}-1\right)
  =
  2^{\lfloor n/2\rfloor}
  +
  2^{\lceil n/2\rceil}
  -2.
\]

For \(n=1,2,\ldots,6\), this upper bound gives
\[
  1,\ 2,\ 4,\ 6,\ 10,\ 14,
\]
which matches the small values listed above.

\begin{conjecture}[Balanced split optimality]
  For every \(n\ge 1\),
  \[
    \tau_2(n)
    =
    2^{\lfloor n/2\rfloor}
    +
    2^{\lceil n/2\rceil}
    -2.
  \]
  Equivalently, the balanced split construction is optimal.
\end{conjecture}

\section{Elementary Lower Bounds}

The first lower bounds are simple but useful.

\begin{observation}[Singletons are forced]
  Every two-foot treadle basis \(\A\) must contain all singleton subsets
  \[
    \{1\},\{2\},\ldots,\{n\}.
  \]
\end{observation}

\begin{proof}
  Suppose \(\{i\}=U\cup V\) with \(U,V\in\A\). Since \(U\) and \(V\) are
  nonempty and their union is \(\{i\}\), both \(U\) and \(V\) must equal
  \(\{i\}\). Hence \(\{i\}\in\A\).
\end{proof}

\begin{proposition}[Counting bound]
  If \(\A\) is a two-foot treadle basis and \(|\A|=m\), then
  \[
    \binom{m+1}{2} \ge 2^n-1.
  \]
  Therefore
  \[
    \tau_2(n)
    \ge
    \left\lceil
      \frac{\sqrt{1+8(2^n-1)}-1}{2}
    \right\rceil.
  \]
\end{proposition}

\begin{proof}
  There are \(m\) choices for a one-treadle union and \(\binom{m}{2}\) choices
  for a union of two distinct treadles. Thus there are at most
  \[
    m+\binom{m}{2}=\binom{m+1}{2}
  \]
  possible unions of at most two members of \(\A\). These unions must include
  all \(2^n-1\) nonempty subsets of \([n]\).
\end{proof}

This counting bound has the right exponential order, namely
\(\tau_2(n)=\Omega(2^{n/2})\), but it does not by itself prove the balanced
split conjecture. The gap is only a constant factor asymptotically, but closing
that gap appears to require using more of the structure of union than just the
number of available pairs.

\section{Equivalent Formulations}

The same problem can be described in several equivalent ways.

\begin{itemize}[leftmargin=2em]
  \item \textbf{Loom language.} Choose the smallest set of treadle tie-ups so
  that every shaft liftplan is available by pressing one or two treadles.

  \item \textbf{Boolean lattice language.} Find the smallest
  \(\A\subseteq\pow([n])\setminus\{\varnothing\}\) such that
  \[
    \pow([n])\setminus\{\varnothing\}
    =
    \{U\cup V : U,V\in\A\}.
  \]

  \item \textbf{Semilattice basis language.} Find the smallest generating set
  for the join-semilattice \((\pow([n])\setminus\{\varnothing\},\cup)\) when
  every element must be generated using at most two basis elements.
\end{itemize}

\section{Questions}

The main open direction is to prove, disprove, or refine the balanced split
conjecture. Some more concrete questions are:

\begin{enumerate}[leftmargin=2em]
  \item Can one prove that every optimal family must have size at least
  \(2^{\lfloor n/2\rfloor}+2^{\lceil n/2\rceil}-2\)?

  \item Are all extremal families essentially balanced split constructions, up
  to relabeling of the shafts?

  \item What changes if the weaver can press at most \(r\) treadles at once
  instead of at most two?

  \item What changes if some treadle tie-ups are physically undesirable, for
  instance because they lift too many shafts at once?
\end{enumerate}

\end{document}
